\section{Score Function 的定义与性质}

\subsection{Score Function 的定义}

Score function（分数函数）在统计学中有精确的定义，它是概率密度函数对数的梯度：

\begin{importantbox}{Score Function 的定义}
给定概率密度函数 $p(x)$，其 score function 定义为：
$$s(x) = \nabla_x \log p(x)$$
其中 $\nabla_x$ 表示关于数据点 $x$ 的梯度（而非关于模型参数的梯度）。
\end{importantbox}

需要特别注意的是，这里的``score''不是通常意义上的``分数''，而是一个在统计学中有特定含义的术语。Score function 是一个向量场——对于每一个数据空间中的点 $x$，它给出一个与 $x$ 同维度的向量，指向使对数概率增大最快的方向。

\begin{knowledgebox}{Score Function 的直观理解}
Score function $\nabla_x \log p(x)$ 可以理解为一个``指导方向''：
\begin{itemize}
    \item 在高概率密度区域附近，score 的大小较小（因为已接近``山顶''）。
    \item 在低概率密度区域，score 指向高概率密度区域的方向。
    \item score 为零的点对应概率密度的极值点。
\end{itemize}
直观地说，score function 告诉我们``应该往哪个方向走，才能到达概率更高的地方''。
\end{knowledgebox}

\subsection{高斯分布的 Score Function}

为了建立直觉，我们来计算高斯分布的 score function。对于条件分布（即前向跳跃分布）：

$$q(x_t \mid x_0) = \mathcal{N}(x_t;\; \sqrt{\bar{\alpha}_t}\, x_0,\; (1-\bar{\alpha}_t) I)$$

其 score function 为：

$$\nabla_{x_t} \log q(x_t \mid x_0) = -\frac{x_t - \sqrt{\bar{\alpha}_t}\, x_0}{1-\bar{\alpha}_t}$$

这个推导很直接：先写出高斯分布的概率密度函数，取对数后只保留与 $x_t$ 相关的项（即二次型项），然后对 $x_t$ 求梯度即可。

\begin{importantbox}{噪声预测与 Score Function 的关系}
回顾 $x_t = \sqrt{\bar{\alpha}_t}\, x_0 + \sqrt{1-\bar{\alpha}_t}\, \epsilon_t$，可以得到：
$$\epsilon_t = \frac{x_t - \sqrt{\bar{\alpha}_t}\, x_0}{\sqrt{1-\bar{\alpha}_t}}$$
因此：
$$\nabla_{x_t} \log q(x_t \mid x_0) = -\frac{\epsilon_t}{\sqrt{1-\bar{\alpha}_t}}$$
即 \textbf{DDPM 中的噪声预测器所预测的 $\epsilon_t$，与条件 score function 之间只差一个缩放因子！} 噪声预测的物理意义就是在学习 score function（up to scale）。
\end{importantbox}

这个发现揭示了 DDPM 噪声预测器的深层含义：表面上我们在``预测噪声''，实际上我们在学习前向跳跃分布的 score function。

\subsection{从条件 Score 到边缘 Score}

上面我们讨论的是\textbf{条件}分布 $q(x_t \mid x_0)$ 的 score function，因为我们知道 $x_0$。但在生成过程中，我们并不知道 $x_0$，我们真正需要的是\textbf{边缘}分布 $q(x_t)$ 的 score function：

$$\nabla_{x_t} \log q(x_t) = ?$$

其中 $q(x_t) = \int q(x_t \mid x_0)\, q(x_0)\, dx_0$。由于 $q(x_0)$（真实数据分布）通常未知且复杂，$q(x_t)$ 一般没有解析形式。

然而，通过数学推导可以证明：

$$\nabla_{x_t} \log q(x_t) = \mathbb{E}_{q(x_0 \mid x_t)}\left[\nabla_{x_t} \log q(x_t \mid x_0)\right]$$

\begin{importantbox}{边缘 Score 是条件 Score 的后验期望}
边缘分布的 score function 等于条件 score function 在后验分布 $q(x_0 \mid x_t)$ 下的期望。这意味着我们训练的噪声预测网络 $\epsilon_\theta(x_t, t)$ 实际上是在学习对所有可能的 $x_0$ 取后验平均的 score function。
\end{importantbox}

\subsection{本章小结}

Score function 定义为对数概率密度的梯度，它编码了关于概率分布的完整信息。高斯分布的 score function 有简洁的闭式解，而 DDPM 中的噪声预测器本质上是在学习前向跳跃分布的 score function（差一个缩放因子）。边缘分布的 score 等于条件 score 在后验下的期望，这为理解扩散模型中神经网络的学习目标提供了严格的数学解释。

